\newcommand{\testNumAfter}{\test{\testNum}}

\begin{center}
\fosfamily\fontseries{bc}\selectfont
\begin{tikzpicture}
  % \draw (0, -1.3) node[opacity=1,inner sep=0] {\logo{0.6}};
    \draw (-5, 0) node 
    {
{\fontsize{30}{50}\selectfont\theYear}};
\draw (-5, -2) node {
{\fontsize{40}{60}\selectfont {\mockAcronym} \tentwelve}};
\draw (-5, -4) node {
{\fontsize{20}{40}\selectfont REMAINS OPEN UNTIL THE DUE DATE}};
\end{tikzpicture}
\rule{\linewidth}{0.2cm}
\LARGE **Administration On An Earlier Date Is Not Even Possible**
\end{center}
\vspace*{5mm}
{
\fontfamily{qpl}\selectfont
\normalfont
\begin{enumerate}
\item All information (Rules and Instructions) needed to administer this exam is contained in the TEACHERS' MANUAL. PLEASE READ THE MANUAL BEFORE \uppercase{\duedate}.
\item  Your PRINCIPAL or VICE-PRINCIPAL must verify on the AMC {\tentwelve} CERTIFICATION FORM (found in the Teachers' Manual) that you followed all rules associated with the conduct of the exam.
\item The Answer Forms must be mailed by trackable mail to the AMC office no later than 24 hours following the exam.
\item \textit{The publication, reproduction or communication of the problems or solutions for this contest during the period when students are eligible to participate seriously jeopardizes the integrity of the results. Dissemination at any time via copier, telephone, email, internet, or media of any type is a violation of the competition rules.}
\end{enumerate}
}

\vspace{4mm}
\begin{center}
\rule{\linewidth}{0.1mm}
\end{center}

    %\begin{center} %\includegraphics[width=14cm]{credits3.PNG}\end{center}
%\begin{center} \includegraphics[width=14cm]{credits2.PNG}\end{center}
\begin{center}
\textit{The {\bfseries\mockName} are brought to you by:}\\[1\baselineskip]
{\bfseries \committee,}
 
 and the {\committeeName}
\end{center}

\newpage
\begin{center}
\fosfamily\fontseries{bc}
\setlength{\fboxrule}{3pt}
\fbox{
\begin{tikzpicture}
    \draw (-15,0) node [anchor=west] {\mockLogo};
    \draw (-13, 0) node[anchor=west] {
{\fontsize{30}{50}\selectfont {\committeeAcronym} {\mockAcronym}}};
\draw (-15, -1.5) node[anchor=west] {
{\fontsize{15}{30}\selectfont\mockName}};
\draw (-15, -2.5) node[anchor=west] {
{\fontsize{15}{30}\selectfont \testNum$^{\text{\testNumAfter}}$ ANNUAL}};
\draw (-15, -4) node[anchor=west] {
{\fontsize{40}{40}\selectfont {\mockAcronym} {\tentwelve}}};
\draw (-15, -5.5) node[anchor=west] {
{\fontsize{15}{30}\selectfont {\mockName} {\tentwelve}}};
\draw [white] (-15,-6) -- (0,-6);
\end{tikzpicture}
}
\end{center}
\begin{center}
\large{\textbf{INSTRUCTIONS}}
\end{center}
\normalfont{
\begin{enumerate}
\item DO NOT OPEN THIS BOOKLET UNTIL THE TIMER STARTS.
\item This is a twenty-five question multiple choice test. Each question is followed by answers marked A, B, C, D and E. Only one of these is correct.
\item Mark your answer clearly, edits in submission will not be accepted.
\item SCORING: You will receive 6 points for each correct answer, 1.5 points for each problem left unanswered, and 0 points for each incorrect answer.
\item No aids are permitted other than scratch paper, graph paper, rulers, compass, protractors, and erasers. No calculators, smartwatches, or computing devices are allowed. No problems on the test will require the use of a calculator.
\item Figures are not necessarily drawn to scale.
\item Before beginning the test, your proctor will ask you to record certain information on the answer form.
\item When your proctor gives the signal, begin working on the problems. You will have 75 minutes to complete the test.
\item When you finish the exam, sign your name in the space provided on the Answer Form.
\end{enumerate}
\vspace{3mm}
\begin{center}
\rule{\linewidth}{0.1mm}
\end{center}
 
\scriptsize{
\textit{The {\committeeName} (\committeeAcronym) reserves the right to re-examine students before deciding whether to grant official status to their scores. Students, regardless of their {\mockAcronym} {\tentwelve} score, will be invited to the {\JIME} ({\JIMEAcronym}).}

}
}
\newpage
